\chapter{General Introduction}

\section{Context}

Billing is one of the largest data sets of a telecom operator. Each call, SMS, data session
and roaming event becomes a billed line, and those lines must be kept for years for customer
care, finance, disputes and legal audits. In many companies this history accumulates in the
same table that the billing application writes to every day. The table grows without limit,
and work that was cheap when it held one year of data becomes slow with six.

This project was carried out within the \Course{} course. Its brief (Project 22) asks for a
data lifecycle strategy that combines partitioning where it is licensed, indexing, archiving
and performance benchmarking, and for proof of its effect on comparable workloads.

\section{Motivation}

Three symptoms of the client's current design drive the work:

\begin{itemize}
  \item \textbf{Reports slow down every month.} A report on last month's revenue reads the
        whole six-year table, because one month is too large a share of the rows for an
        index to help.
  \item \textbf{Retention is expensive.} Removing a month means \code{DELETE}-ing hundreds of
        thousands of rows. Every row and every index entry generates undo and redo, and the
        table does not even shrink afterwards.
  \item \textbf{Retention is manual and risky.} Nothing prevents someone from deleting data
        that is still under legal retention, and nothing proves afterwards that no billing
        record was lost or changed.
\end{itemize}

\section{Problem statement}

\emph{How can the billing history be organised so that frequent reports read only the period
they need, old data moves to cheaper storage and is removed at the end of its retention with
near-zero cost, while guaranteeing, and proving, that no billing record is lost or altered?}

\section{Objectives}

\begin{enumerate}
  \item Generate a representative, reproducible billing history (tens of millions of rows).
  \item Measure the performance of the current design on a defined workload: the
        \emph{baseline}.
  \item Design a partitioned schema and a lifecycle policy (hot, warm, archive, purge) suited
        to the client's retention rules and licences.
  \item Implement the lifecycle as automated, guarded, logged and recoverable database code.
  \item Measure query, maintenance and storage effects on the same workload, and verify data
        integrity objectively.
  \item Document the licensing restrictions, trade-offs and limits of the solution.
\end{enumerate}

\section{Structure of the report}

Chapter~\ref{ch:client} presents the client. Chapters~\ref{ch:requirements}
and~\ref{ch:technology} analyse the requirements and compare the possible technologies.
Chapter~\ref{ch:architecture} describes the architecture, and
Chapters~\ref{ch:environment} and~\ref{ch:implementation} the environment and the
implementation. Chapter~\ref{ch:usage} shows practical scenarios. Chapters~\ref{ch:testing}
and~\ref{ch:performance} present the tests and the measured results.
Chapters~\ref{ch:challenges} to~\ref{ch:conclusion} discuss the difficulties, the limits and
the outcome. All the code, test outputs and benchmark data are in the project repository
(\url{\RepoURL}).
